%% defense_eval.tex -- submission version

\section{RQ3: Mitigation Boundary and Legitimate Recovery}
\label{sec:defense-eval}

RQ3 separates mechanism screens from native-framework exploitability. The
completed defense-aware and legitimate-recovery results are descriptive at
their sample sizes; earlier naive-baseline and cross-model protocols are not
included as quantitative claims because their raw artifacts and snapshot
provenance are incomplete.

\subsection{Defense-aware Attacker}
\label{sec:defense-aware}

The defense-aware screen uses five payload classes, 20 trials per class and
configuration, and GPT-5.6 Luna. The undefended arm succeeds in
\DefenseUndefendedSuccesses{}/\DefenseTrialsPerConfig{} trials
(\DefenseUndefendedASR{}); the defended arm succeeds in
\DefenseDefendedSuccesses{}/\DefenseTrialsPerConfig{} trials
(\DefenseDefendedASR{}). Fisher's exact test gives
$p=\DefenseFisherP{}$. The Wilson 95\% intervals are
[\DefenseUndefendedCILow{}, \DefenseUndefendedCIHigh{}] and
[\DefenseDefendedCILow{}, \DefenseDefendedCIHigh{}], respectively. No trial
was blocked by an API error or a forward-path check.

The point estimate is consistent with a modest reduction, but the contrast is
not statistically established at the completed sample size. The power
analysis estimates approximately 300 trials per pooled arm for 80\% power for
this observed difference under the planning approximation. We therefore report
the result as supplementary evidence, not as a general defense-effectiveness
claim.

\subsection{Legitimate Recovery}
\label{sec:defense-cost}

We generated \LegalRecoveryTasks{} deterministic tasks across arithmetic,
assertion failure, compile error, external-content error, format error,
permission denial, and tool timeout. GPT-5.6 Luna produced
\LegalRecoverySuccesses{}/\LegalRecoveryTasks{} quality-passing repairs and
triggered zero false positives. This supports the feasibility of benign
recovery under the tested task generator, while the task set remains a
screening sample rather than a population-quality estimate.

\paragraph{RQ3 summary.}
The completed screen supports two cautious statements: defense-aware recovery
remains vulnerable in the tested configuration, and legitimate recovery is
feasible with no observed false positives. It does not support a claim that
\trustorigin{} significantly reduces attack success across models or
frameworks.
